\chapter{Introducción y Objetivos}

El osciloscopio de almacenamiento digital (DSO) con módulo matemático FFT que se emplea para efectuar este trabajo
práctico permite observar y analizar una misma señal tanto en el dominio del tiempo como en el de la frecuencia.
En esta actividad de laboratorio se busca poner en evidencia las ventajas que supone el estudio de una señal en el
dominio frecuencial, particularmente cuando la misma contiene una cantidad apreciable de armónicos o ruido superpuesto.

Dado que en este trabajo se utiliza un osciloscopio digital, es necesaria una previa lectura del manual del instrumento
a fin de familiarizarse con cada uno de los menús que permiten controlar sus funciones: adquisición, modos de visualización,
cursores y menú matemático.
La mayoría de los osciloscopios digitales poseen una función de Autoset que ajusta automáticamente la base de tiempos y
la sensibilidad vertical a valores adecuados, lo cual facilita la operación inicial del instrumento. No obstante, en las
experiencias se procede mediante el ajuste manual riguroso de cada parámetro para dominar acabadamente la operación del
equipo.

\section{Objetivos}
En conformidad con la consigna de la cátedra, los objetivos fijados para la práctica son:
\begin{itemize}
    \item Familiarizar al alumno con el empleo de osciloscopios digitales y su módulo matemático FFT.
    \item Estudiar distintas señales en el dominio de la frecuencia y utilizar las utilidades de los menús de medidas
    que suelen poseer estos instrumentos.
    \item Dar el resultado de las mediciones acompañadas del grado de incertidumbre.
\end{itemize}

\section{Consideraciones Previas y Criterio de Incertidumbre}
Antes de comenzar con los ensayos, es indispensable identificar claramente cuáles son las placas o dispositivos bajo
prueba que se utilizarán y tomar debida nota de los mismos. Asimismo, todo proceso de medición requiere contar con una
estimación previa de los valores esperados para corroborar la validez del método y los instrumentos seleccionados.

En cumplimiento de los objetivos fijados, las mediciones cuantitativas se presentan acompañadas de su correspondiente
estimación de incertidumbre expandida para un factor de cobertura $k = 2$, representativo de un nivel de confianza del
$95{,}45\%$.
